% This document was generated by an LLM.

\documentclass{workbook}

\title{Inequalities in the Real Numbers\\
\large A Self-Check and Proof Workbook}
\author{}
\date{}

\begin{document}
\maketitle

\noindent
This workbook develops inequalities from the ordered-field axioms through
absolute values, triangle inequalities, and quantitative estimates. It
contains questions and proof requests, not solutions. The problems are roughly
ordered by logical dependency rather than by computational difficulty.

\medskip\noindent
\textbf{Working convention.} Unless a problem says otherwise, all variables
denote real numbers. You may use the field axioms, the axioms of a total order,
and compatibility of that order with addition and multiplication. When a
problem asks you to prove a derived rule, do not silently use that same rule
in its proof. You may cite earlier results once you have proved them.
Distinguish the hypotheses $x>0$, $x\geq 0$, and $x\neq 0$ carefully.

\medskip\noindent
\textbf{How to work.} For a recall question, state the result precisely,
including its assumptions. For a proof, justify each nontrivial implication.
For a ``decide'' question, either prove the statement or give a concrete
counterexample. In longer arguments, identify which previously proved lemmas
you use.

\section{Axioms, definitions, and the language of order}

\begin{problem}{Recall.}
State the additive and multiplicative parts of the field axioms. Which
operations are required to have inverses, and for which elements?
\end{problem}

\begin{problem}{Recall.}
What does it mean for $\leq$ to be a total (linear) order? State
reflexivity, antisymmetry, transitivity, and totality.
\end{problem}

\begin{problem}{Recall.}
Define $a<b$ in terms of $a\leq b$ and equality. What is the trichotomy
property for two real numbers? How are $>$ and $\geq$ defined?
\end{problem}

\begin{problem}{Recall.}
State the compatibility of order with addition. Give one standard axiom
formulating compatibility with multiplication of nonnegative elements.
Explain why multiplication by a negative element is \emph{not} part of the
positive-multiplication rule.
\end{problem}

\begin{problem}{Concept.}
Which axioms distinguish the real numbers from an arbitrary ordered field?
Do you expect to need completeness to prove elementary absolute-value
inequalities? Explain briefly.
\end{problem}

\begin{problem}{Recall.}
Distinguish an axiom, a definition, a lemma, and a theorem. Why is it
legitimate to invoke a proved lemma without repeating its entire proof?
\end{problem}

\begin{problem}{Prove.}
Show that $a\leq b$ if and only if $0\leq b-a$. Use compatibility with
addition and the additive inverse axioms only.
\end{problem}

\begin{problem}{Prove.}
Show that $a\leq b$ if and only if $a+c\leq b+c$. Justify both directions;
do not confuse the result with the one-way order axiom.
\end{problem}

\begin{problem}{Prove.}
If $a\leq b$ and $c\leq d$, show that $a+c\leq b+d$. Write a two-step
inequality chain and identify where transitivity is used.
\end{problem}

\begin{problem}{Prove.}
If $a<b$ and $c\leq d$, prove $a+c<b+d$. Check carefully where a strict
inequality is preserved.
\end{problem}

\section{Deriving rules for manipulating inequalities}

\begin{problem}{Prove.}
Starting from $a\leq b$, derive $-b\leq -a$ by adding suitable terms to
both sides. Do not use multiplication by $-1$ as an already known rule.
\end{problem}

\begin{problem}{Prove.}
Using only field axioms, establish $-(-a)=a$ and
$-(a+b)=(-a)+(-b)$. Explain the uniqueness of the additive inverse.
\end{problem}

\begin{problem}{Prove.}
Establish $(-1)a=-a$. Then explain why ``multiply by $-1$'' and ``take
additive inverses'' describe the same operation, even though their order
properties must still be proved.
\end{problem}

\begin{problem}{Prove.}
Show that $a<b$ implies $-b<-a$. Can you obtain this directly from the
non-strict version and the fact that $a\neq b$?
\end{problem}

\begin{problem}{Prove.}
From $a\leq b$ and $c\leq d$, derive $a-d\leq b-c$.
Explain why the superficially similar assertion $a-c\leq b-d$ need not hold.
\end{problem}

\begin{problem}{Prove.}
Assume $0\leq c$ and $a\leq b$. Derive $ac\leq bc$ from the ordered-field
axioms, perhaps by examining $(b-a)c$.
\end{problem}

\begin{problem}{Prove.}
Assume $c<0$ and $a\leq b$. Show that $ac\geq bc$.
Make explicit why $-c>0$, and separate the use of the multiplication rule
from the use of additive inverses.
\end{problem}

\begin{problem}{Prove.}
If $a<b$ and $c>0$, show that $ac<bc$. Which property of a field ensures
that the product of two nonzero elements is nonzero?
\end{problem}

\begin{problem}{Prove.}
If $a\leq b$ and $c\leq d$, with $a,b,c,d\geq0$, prove $ac\leq bd$.
Is it enough to multiply the two original inequalities without discussing
signs?
\end{problem}

\begin{problem}{Recall and prove.}
Formulate the rule for dividing an inequality by a positive number, and
the corresponding rule for a negative number. Derive both rules from
multiplication rather than taking them as new axioms.
\end{problem}

\begin{problem}{Decide.}
If $a\leq b$ and $c\leq d$, is it necessarily true that $a-c\leq b-d$?
Either prove the claim or find a numerical counterexample.
\end{problem}

\begin{problem}{Decide.}
If $a<b$ and $c<d$, must $ac<bd$? If not, find a counterexample and
formulate useful sign assumptions under which a product comparison
\emph{does} follow.
\end{problem}

\section{Signs, squares, and reciprocals}

\begin{problem}{Prove.}
Show that the product of two strictly positive numbers is strictly
positive. Use the nonnegative multiplication axiom and the absence of
zero divisors.
\end{problem}

\begin{problem}{Prove.}
Prove $a^2\geq0$ for every real $a$, and $a^2>0$ whenever $a\neq0$.
Deal explicitly with the possibility $a<0$.
\end{problem}

\begin{problem}{Prove.}
Use a preceding result and $1\neq0$ to prove $0<1$.
Deduce $-1<0$.
\end{problem}

\begin{problem}{Prove.}
If $0\leq a\leq b$, show that $a^2\leq b^2$.
Give one proof using multiplication of inequalities and another using
$b^2-a^2=(b-a)(b+a)$.
\end{problem}

\begin{problem}{Decide.}
Is $a\leq b\Rightarrow a^2\leq b^2$ valid without sign assumptions?
If not, state a corrected version and prove it.
\end{problem}

\begin{problem}{Prove.}
Show that $(a-b)^2\geq0$ implies $2ab\leq a^2+b^2$.
Characterize the equality case.
\end{problem}

\begin{problem}{Prove.}
If $a>0$, prove $1/a>0$. Justify your argument without assuming in
advance that reciprocals preserve positivity.
\end{problem}

\begin{problem}{Prove.}
If $0<a<b$, prove $0<1/b<1/a$.
Which of these inequalities become false or require modification when
$a$ and $b$ are negative?
\end{problem}

\begin{problem}{Prove.}
For nonzero $a,b$ satisfying $a<b$ and $ab>0$, show $1/a>1/b$.
Use the identity $1/a-1/b=(b-a)/(ab)$, justifying its sign.
\end{problem}

\begin{problem}{Decide.}
Does $a^2=b^2$ imply $a=b$? State the exact conclusion that follows,
and specify an additional hypothesis that makes $a=b$ valid.
\end{problem}

\begin{problem}{Prove.}
For $a,b\geq0$, show that $a^2\leq b^2$ if and only if $a\leq b$.
Do not invoke square roots unless you also justify their use.
\end{problem}

\section{Absolute value from the definition}

\begin{problem}{Recall.}
Give the piecewise definition of $|x|$. What does this definition say
when $x=0$? State the range of the absolute-value function.
\end{problem}

\begin{problem}{Prove.}
Show that $|x|\geq0$ and that $|x|=0$ if and only if $x=0$.
\end{problem}

\begin{problem}{Prove.}
Derive $x\leq|x|$ and $-x\leq|x|$ from the definition, considering cases
if necessary. Deduce $-|x|\leq x\leq|x|$.
\end{problem}

\begin{problem}{Prove.}
Show that $|-x|=|x|$ in two ways: first from the piecewise definition,
then from $|x|^2=x^2$. In the second proof, explain why equality of
squares is sufficient \emph{here}.
\end{problem}

\begin{problem}{Prove.}
Establish $|x|^2=x^2$ and $\bigl||x|\bigr|=|x|$.
\end{problem}

\begin{problem}{Prove.}
For $M\geq0$, establish the equivalence
$|x|\leq M\iff -M\leq x\leq M$.
Prove both implications, and check $M=0$ explicitly.
\end{problem}

\begin{problem}{Investigate.}
Does the preceding equivalence remain logically true when $M<0$?
Analyze whether either side can hold. Distinguish an equivalence from a
claim that the conditions are satisfiable.
\end{problem}

\begin{problem}{Prove.}
For $M>0$, prove $|x|<M\iff -M<x<M$.
What happens when $M=0$?
\end{problem}

\begin{problem}{Prove.}
For $M\geq0$, show that $|x|\geq M$ is equivalent to
$x\leq -M$ or $x\geq M$. Explain the role of the word ``or.''
\end{problem}

\begin{problem}{Prove.}
Show that $|ab|=|a|\,|b|$ for all real $a,b$.
Try one proof by cases and a different proof by comparing squares of
nonnegative quantities.
\end{problem}

\begin{problem}{Prove.}
For $b\neq0$, prove $|a/b|=|a|/|b|$.
Explain why $|b|>0$.
\end{problem}

\begin{problem}{Prove.}
Establish $|a-b|=|b-a|$ without treating separately the cases $a>b$,
$a=b$, and $a<b$.
\end{problem}

\begin{problem}{Prove.}
Show that $|x|=\max\{x,-x\}$.
Then prove
$\max\{a,b\}=(a+b+|a-b|)/2$ and
$\min\{a,b\}=(a+b-|a-b|)/2$.
\end{problem}

\section{The triangle inequality and its relatives}

\begin{problem}{Prove.}
Prove $|a+b|\leq|a|+|b|$ directly from
$-|x|\leq x\leq|x|$ and the addition compatibility axiom.
Do not invoke any previously known triangle inequality.
\end{problem}

\begin{problem}{Prove.}
Prove $|a-b|\leq|a|+|b|$ directly from order bounds, without invoking
either the triangle inequality or $|a-b|=|b-a|$.
\end{problem}

\begin{problem}{Prove again.}
Give a one-line proof of $|a-b|\leq|a|+|b|$ using the triangle inequality
and the identity $|-b|=|b|$. Compare its logical dependencies with those
of the preceding proof.
\end{problem}

\begin{problem}{Prove.}
Derive $|a|\leq|a-b|+|b|$ by rewriting $a$ as a sum.
\end{problem}

\begin{problem}{Prove.}
Establish the \emph{reverse triangle inequality}
$\bigl||a|-|b|\bigr|\leq|a-b|$.
First obtain a one-sided bound; then exchange $a$ and $b$.
\end{problem}

\begin{problem}{Prove.}
Derive $\bigl||a|-|b|\bigr|\leq|a+b|$.
Be explicit about which sign change turns this into an application of
the preceding result.
\end{problem}

\begin{problem}{Prove.}
Combine your results to establish
$\bigl||a|-|b|\bigr|\leq|a-b|\leq|a|+|b|$.
Find numerical examples where each inequality is strict, and others
where equality holds.
\end{problem}

\begin{problem}{Prove.}
Show that $|a+b+c|\leq|a|+|b|+|c|$.
Extend the assertion to $n$ real summands by mathematical induction.
\end{problem}

\begin{problem}{Characterize equality.}
For real $a,b$, prove that $|a+b|=|a|+|b|$ if and only if $ab\geq0$.
Be sure to include zero cases.
\end{problem}

\begin{problem}{Characterize equality.}
Determine exactly when $|a-b|=|a|+|b|$.
Give an algebraic or case-based proof of your condition.
\end{problem}

\begin{problem}{Prove.}
Establish $2|a|\leq|a+b|+|a-b|$.
You may express $2a$ as a sum involving $a+b$ and $a-b$.
\end{problem}

\begin{problem}{Challenge.}
Prove the real-number identity
$|a+b|+|a-b|=2\max\{|a|,|b|\}$.
Try first to see why this is stronger than the preceding estimate.
\end{problem}

\begin{problem}{Concept and proof.}
Define $d(a,b)=|a-b|$ on $\mathbb R$. Verify nonnegativity,
$d(a,b)=0\iff a=b$, symmetry, and the metric triangle inequality.
Which parts have you already proved in another form?
\end{problem}

\section{Working with bounds and error estimates}

\begin{problem}{Prove.}
For $r>0$, show that $|x-c|<r$ is equivalent to $c-r<x<c+r$.
Write the version with weak inequalities as well.
\end{problem}

\begin{problem}{Prove.}
If $|x-a|\leq r$ and $|x-b|\leq s$, where $r,s\geq0$, prove
$|a-b|\leq r+s$.
\end{problem}

\begin{problem}{Prove.}
If $|x-a|<r$ and $|a-b|<s$, with $r,s>0$, prove
$|x-b|<r+s$. Identify where strictness enters.
\end{problem}

\begin{problem}{Prove.}
If $\varepsilon>0$, $|x|<\varepsilon/2$, and
$|y|<\varepsilon/2$, prove $|x+y|<\varepsilon$.
Generalize to bounds $|x|<\alpha$ and $|y|<\beta$.
\end{problem}

\begin{problem}{Prove.}
If $|a|\leq A$ and $|b|\leq B$, with $A,B\geq0$, prove
$|ab|\leq AB$. Name every sign assumption used.
\end{problem}

\begin{problem}{Prove.}
Use factorization to establish
$|a^2-b^2|=|a-b|\,|a+b|$.
Deduce $|a^2-b^2|\leq 2R|a-b|$ whenever
$|a|,|b|\leq R$, $R\geq0$.
\end{problem}

\begin{problem}{Prove.}
For $a,b\neq0$, obtain a formula for $|1/a-1/b|$ in terms of
$|a-b|$, $|a|$, and $|b|$. Deduce an upper bound when
$|a|,|b|\geq m>0$.
\end{problem}

\begin{problem}{Prove.}
Establish
$|ab-cd|\leq |a|\,|b-d|+|d|\,|a-c|$.
Find the useful intermediate expression obtained by adding and
subtracting the same product.
\end{problem}

\begin{problem}{Challenge.}
Let $n\geq2$ be an integer and $|x|,|y|\leq R$ for some $R>0$.
Prove $|x^n-y^n|\leq nR^{n-1}|x-y|$ using factorization and the
finite-sum triangle inequality.
\end{problem}

\begin{problem}{Prove.}
If $|x|\leq\varepsilon$ for every real $\varepsilon>0$, show that $x=0$.
Explain why this result is a statement about the order of real numbers,
not a theorem about limits.
\end{problem}

\section{Proof audits and counterexamples}

\begin{problem}{Audit.}
A purported proof says: ``$a\leq|a|$ and $b\leq|b|$, hence
$a-b\leq|a|-|b|$.'' Identify the invalid inference and repair the
argument if a useful inequality was intended.
\end{problem}

\begin{problem}{Audit.}
A purported proof says: ``$u^2=v^2$, therefore $u=v$.''
What is missing? Give a counterexample and a sufficient extra hypothesis.
\end{problem}

\begin{problem}{Audit.}
A purported proof says: ``Since $a<b$, multiplying by $c$ gives
$ac<bc$.'' List the possibilities for the sign of $c$ and give the
correct conclusion in each case.
\end{problem}

\begin{problem}{Audit.}
A purported proof of $|x+y|\leq|x|+|y|$ begins by applying the triangle
inequality to $x$ and $y$. Explain the logical problem, even though the
result stated is true.
\end{problem}

\begin{problem}{Decide.}
Does $|a|\leq|b|$ imply $a\leq b$? Does $a\leq b$ imply $|a|\leq|b|$?
Treat the two implications separately.
\end{problem}

\begin{problem}{Decide.}
If $|a+b|=|a|+|b|$, must $a,b\geq0$?
If not, give the exact weaker condition.
\end{problem}

\begin{problem}{Decide.}
If $|x|<\varepsilon$ with $\varepsilon>0$, must $x<\varepsilon$?
Does the converse hold? Supply a proof or counterexample for each direction.
\end{problem}

\begin{problem}{Decide.}
For real $a,b,c$, is
$|a-c|\leq |a-b|+|b-c|$ always true?
Provide a proof that does not appeal to geometric intuition.
\end{problem}

\begin{problem}{Audit.}
Someone claims that $-M\leq x\leq M$ can be converted to
$|x|\leq M$ only when $M>0$. Analyze separately $M=0$ and $M<0$.
Is the strict positivity assumption essential to the \emph{equivalence}?
\end{problem}

\begin{problem}{Synthesis.}
Write a self-contained proof of
$\bigl||a|-|b|\bigr|\leq|a-b|\leq|a|+|b|$.
Before writing the proof, list its dependencies as a small directed
graph of axioms, definitions, and lemmas. Then remove any dependency
that is unnecessary.
\end{problem}

\section*{Final self-check}

Without consulting notes, explain why an order-preserving operation needs
sign hypotheses when multiplication is involved but not when addition is
involved. State the assumptions behind each rule you used. Reconstruct a
proof of the ordinary and reverse triangle inequalities. Finally, give a
counterexample to at least three plausible but false inequality rules.

\end{document}
